\documentclass[aspectratio=169,xcolor=dvipsnames]{beamer}

\usetheme{SimpleDarkBlue}

% ------------------------------------------------
% Codificação e idioma
% ------------------------------------------------

\usepackage[T1]{fontenc}
\usepackage[utf8]{inputenc}
\usepackage[brazilian]{babel}
\usepackage{lmodern}

% ------------------------------------------------
% Matemática
% ------------------------------------------------

\usepackage{amsmath}
\usepackage{amsthm}
\usepackage{amssymb}
\usepackage{mathtools}

% ------------------------------------------------
% Outros
% ------------------------------------------------

\usepackage{hyperref}
\usepackage{graphicx}
\usepackage{booktabs}
\usepackage{listings}
\usepackage{xcolor}

% ------------------------------------------------
% Configurações
% ------------------------------------------------

\setbeamertemplate{caption}{\insertcaption}

\definecolor{codebackground}{RGB}{245,245,245}
\definecolor{codeblue}{RGB}{30,80,150}
\definecolor{codegreen}{RGB}{30,120,70}
\definecolor{codegray}{RGB}{100,100,100}

\lstset{
    basicstyle=\ttfamily\small,
    backgroundcolor=\color{codebackground},
    frame=single,
    rulecolor=\color{gray!30},
    keywordstyle=\color{codeblue},
    commentstyle=\color{codegreen},
    breaklines=true,
    showstringspaces=false,
    columns=fullflexible,
    keepspaces=true
}

\usepackage{listings}
\lstset{
    basicstyle=\ttfamily\small,
    backgroundcolor=\color{codebackground},
    frame=single,
    rulecolor=\color{gray!30},
    keywordstyle=\color{codeblue},
    commentstyle=\color{codegreen},
    breaklines=true,
    showstringspaces=false,
    columns=fullflexible,
    keepspaces=true,
    inputencoding=utf8,
    extendedchars=false,
    literate={nbsp}{{\ }}1
}

% ------------------------------------------------
% Informações
% ------------------------------------------------

\title{Fórmulas Matemáticas}
\subtitle{Sintaxe, Estruturas e Notação Matemática}

\author{Rodrigo Martins \and Pedro Motta}

\institute{
    Minicurso de Introdução ao LaTeX\\
    Universidade Federal Rural do Rio de Janeiro
}

\date{}

% =================================================
% DOCUMENTO
% =================================================

\begin{document}

% ------------------------------------------------
% Capa
% ------------------------------------------------

\begin{frame}
    \titlepage
\end{frame}

% ------------------------------------------------
% Objetivos
% ------------------------------------------------

\section{Objetivos}

\begin{frame}{O que vamos aprender?}

    \begin{itemize}
        \item Sintaxe matemática do LaTeX;
        \item Funções e operadores;
        \item Matrizes e vetores;
        \item Sistemas de equações;
        \item Notações matemáticas avançadas;
        \item Combinação de diferentes estruturas.
    \end{itemize}

\end{frame}

\begin{frame}{A ideia central}

    \begin{block}{LaTeX não desenha a fórmula}
        Nós descrevemos a \textbf{estrutura} da expressão matemática.
    \end{block}

    \vspace{0.5cm}

    \[
        \frac{x^2+1}{x-1}
    \]

    \vspace{0.3cm}

    \begin{center}
        \texttt{\textbackslash frac\{x\^2+1\}\{x-1\}}
    \end{center}

\end{frame}

% =================================================
% MODO MATEMÁTICO
% =================================================

\section{Modo Matemático}

\begin{frame}[fragile]{Entrando no modo matemático}

    \begin{block}{Inline}
        Matemática dentro do texto.
    \end{block}

\begin{lstlisting}
A equacao $x^2+1=0$ possui raizes complexas.
\end{lstlisting}

    \begin{block}{Destacado}
        Matemática em uma linha própria.
    \end{block}

\begin{lstlisting}
\[
x^2+1=0
\]
\end{lstlisting}

\end{frame}

\begin{frame}{Inline ou destacado?}

    \textbf{Inline:}

    \[
        \text{A equação } x^2+1=0
    \]

    \vspace{0.4cm}

    \textbf{Destacado:}

    \[
        x^2+1=0
    \]

    \vspace{0.3cm}

    \begin{center}
        Escolha de acordo com o papel da expressão no texto.
    \end{center}

\end{frame}

\begin{frame}[fragile]{Equações numeradas}

    Para equações que precisam de numeração:

\begin{lstlisting}
\begin{equation}
E = mc^2
\end{equation}
\end{lstlisting}

    \begin{equation}
        E=mc^2
    \end{equation}

    \begin{center}
        Útil quando a equação será discutida ou referenciada posteriormente.
    \end{center}

\end{frame}

% =================================================
% SINTAXE
% =================================================

\section{Sintaxe Matemática}

\begin{frame}[fragile]{A estrutura dos comandos}

    A forma mais comum:

\begin{lstlisting}
\comando{argumento}
\end{lstlisting}

    Exemplo:

\begin{lstlisting}
\frac{a}{b}
\end{lstlisting}

    \[
        \frac{a}{b}
    \]

    \begin{alertblock}{Ideia}
        O comando determina \textbf{o que fazer}; as chaves determinam \textbf{com o quê}.
    \end{alertblock}

\end{frame}

\begin{frame}[fragile]{Sobrescritos e subscritos}

    \begin{columns}

        \column{0.48\textwidth}

        \textbf{Sobrescrito}

\begin{lstlisting}
x^2
\end{lstlisting}

        \[
            x^2
        \]

        \column{0.48\textwidth}

        \textbf{Subscrito}

\begin{lstlisting}
x_i
\end{lstlisting}

        \[
            x_i
        \]

    \end{columns}

    \vspace{0.5cm}

    Quando há mais de um caractere:

\begin{lstlisting}
x^{n+1}       a_{ij}
\end{lstlisting}

\[
    x^{n+1}
    \qquad
    a_{ij}
\]

\end{frame}

% =================================================
% SIMBOLOGIA
% =================================================

\begin{frame}[fragile]{Operações básicas}

\begin{lstlisting}
a+b
a-b
a\cdot b
\frac{a}{b}
x^2
\sqrt{x}
\end{lstlisting}

    \[
        a+b
        \qquad
        a-b
        \qquad
        a\cdot b
    \]

    \[
        \frac{a}{b}
        \qquad
        x^2
        \qquad
        \sqrt{x}
    \]

\end{frame}

\begin{frame}[fragile]{Frações}

    Estrutura:

\begin{lstlisting}
\frac{numerador}{denominador}
\end{lstlisting}

    Exemplo:

\begin{lstlisting}
\frac{x+1}{x-1}
\end{lstlisting}

    \[
        \frac{x+1}{x-1}
    \]

    \vspace{0.3cm}

    Estruturas podem ser combinadas:

\begin{lstlisting}
\frac{x^2+1}{\sqrt{x+1}}
\end{lstlisting}

\end{frame}

\begin{frame}[fragile]{Raízes}

    \textbf{Raiz quadrada}

\begin{lstlisting}
\sqrt{x}
\end{lstlisting}

    \[
        \sqrt{x}
    \]

    \textbf{Raiz de índice $n$}

\begin{lstlisting}
\sqrt[n]{x}
\end{lstlisting}

    \[
        \sqrt[n]{x}
    \]

\end{frame}

\begin{frame}[fragile]{Delimitadores}

    Delimitadores simples:

\begin{lstlisting}
(x+1)
[x+1]
\end{lstlisting}

    Para expressões grandes:

\begin{lstlisting}
\left( \frac{x+1}{x-1} \right)
\end{lstlisting}

    \[
        \left(
        \frac{x+1}{x-1}
        \right)
    \]

    \begin{center}
        \texttt{\textbackslash left} e \texttt{\textbackslash right}
        ajustam o tamanho dos delimitadores.
    \end{center}

\end{frame}

\begin{frame}[fragile]{Símbolos gregos}

    Alguns dos mais utilizados:

\begin{lstlisting}
\alpha   \beta   \gamma
\theta   \lambda \mu
\pi      \sigma  \omega
\end{lstlisting}

    \[
        \alpha \quad \beta \quad \gamma
    \]

    \[
        \theta \quad \lambda \quad \mu
    \]

    \[
        \pi \quad \sigma \quad \omega
    \]

\end{frame}

% =================================================
% FUNÇÕES
% =================================================

\section{Funções e Operadores}

\begin{frame}[fragile]{Funções}

    Uma função pode ser escrita diretamente:

\begin{lstlisting}
f(x)=x^2+2x+1
\end{lstlisting}

    \[
        f(x)=x^2+2x+1
    \]

    Estruturas podem ser combinadas:

\begin{lstlisting}
f(x)=\frac{x^2+1}{x-1}
\end{lstlisting}

    \[
        f(x)=\frac{x^2+1}{x-1}
    \]

\end{frame}

\begin{frame}[fragile]{Operadores matemáticos}

    Use comandos próprios para operadores:

\begin{lstlisting}
\sin(x)    \cos(x)    \tan(x)
\log(x)    \ln(x)     \exp(x)
\end{lstlisting}

    \[
        \sin(x)
        \qquad
        \cos(x)
        \qquad
        \tan(x)
    \]

    \[
        \log(x)
        \qquad
        \ln(x)
        \qquad
        \exp(x)
    \]

\end{frame}

\begin{frame}[fragile]{Por que usar \texttt{\textbackslash sin}?}

    \textbf{Evite:}

\begin{lstlisting}
sin(x)
\end{lstlisting}

    \textbf{Prefira:}

\begin{lstlisting}
\sin(x)
\end{lstlisting}

    \[
        \sin(x)
    \]

    \begin{alertblock}{Boa prática}
        Utilize os comandos semânticos fornecidos pelo LaTeX.
    \end{alertblock}

\end{frame}

% =================================================
% EQUAÇÕES LONGAS
% =================================================

\section{Matrizes e Vetores}

\begin{frame}[fragile]{Matrizes}

    Estrutura básica:

\begin{lstlisting}
\begin{matrix}
a & b\\
c & d
\end{matrix}
\end{lstlisting}

    \[
        \begin{matrix}
            a & b\\
            c & d
        \end{matrix}
    \]

    \begin{itemize}
        \item \texttt{\&} separa colunas;
        \item \texttt{\textbackslash\textbackslash} cria uma nova linha.
    \end{itemize}

\end{frame}

\begin{frame}[fragile]{Tipos de matrizes}

    \textbf{Parênteses}

\begin{lstlisting}
\begin{pmatrix}
a & b\\
c & d
\end{pmatrix}
\end{lstlisting}

    \[
        \begin{pmatrix}
            a & b\\
            c & d
        \end{pmatrix}
    \]

    \textbf{Colchetes}

\begin{lstlisting}
\begin{bmatrix}
a & b\\
c & d
\end{bmatrix}
\end{lstlisting}

\end{frame}

\begin{frame}[fragile]{Vetores}

    Vetor coluna:

\begin{lstlisting}
\begin{pmatrix}
x\\
y\\
z
\end{pmatrix}
\end{lstlisting}

    \[
        \begin{pmatrix}
            x\\
            y\\
            z
        \end{pmatrix}
    \]

    \vspace{0.2cm}

    A mesma estrutura pode representar uma matriz $3\times1$.

\end{frame}

\begin{frame}[fragile]{Matrizes com índices}

\begin{lstlisting}
A =
\begin{pmatrix}
a_{11} & a_{12}\\
a_{21} & a_{22}
\end{pmatrix}
\end{lstlisting}

    \[
        A=
        \begin{pmatrix}
            a_{11} & a_{12}\\
            a_{21} & a_{22}
        \end{pmatrix}
    \]

    \begin{center}
        Uma estrutura pode conter outras estruturas.
    \end{center}

\end{frame}

\section{Sistemas de Equações}

\begin{frame}[fragile]{Equações alinhadas}

\begin{lstlisting}
\begin{align}
x+y &= 10\\
x-y &= 2
\end{align}
\end{lstlisting}

    \[
        \begin{aligned}
            x+y &= 10\\
            x-y &= 2
        \end{aligned}
    \]

    \begin{itemize}
        \item \texttt{\&} define o ponto de alinhamento;
        \item \texttt{\textbackslash\textbackslash} quebra a linha.
    \end{itemize}

\end{frame}

\begin{frame}[fragile]{Sistemas de equações}

\begin{lstlisting}
\begin{cases}
x+y=10\\
x-y=2
\end{cases}
\end{lstlisting}

    \[
        \begin{cases}
            x+y=10\\
            x-y=2
        \end{cases}
    \]

    \begin{center}
        O ambiente \texttt{cases} cria a estrutura do sistema.
    \end{center}

\end{frame}

\begin{frame}[fragile]{Funções definidas por partes}

\begin{lstlisting}
f(x)=
\begin{cases}
x^2, & x\geq 0\\
-x, & x<0
\end{cases}
\end{lstlisting}

    \[
        f(x)=
        \begin{cases}
            x^2, & x\geq0\\
            -x, & x<0
        \end{cases}
    \]

\end{frame}




\begin{frame}[fragile]{Limites}

    Estrutura:

\begin{lstlisting}
\lim_{x\to a} f(x)
\end{lstlisting}

    Exemplo:

\begin{lstlisting}
\lim_{x\to 0}\frac{\sin x}{x}=1
\end{lstlisting}

    \[
        \lim_{x\to 0}
        \frac{\sin x}{x}
        =1
    \]

\end{frame}

\begin{frame}[fragile]{Somatórios}

    Estrutura:

\begin{lstlisting}
\sum_{i=1}^{n} x_i
\end{lstlisting}

    \[
        \sum_{i=1}^{n}x_i
    \]

    Outro exemplo:

\begin{lstlisting}
\sum_{i=1}^{n} i^2
\end{lstlisting}

    \[
        \sum_{i=1}^{n}i^2
    \]

\end{frame}

\begin{frame}[fragile]{Produtórios}

    A estrutura é semelhante ao somatório:

\begin{lstlisting}
\prod_{i=1}^{n} x_i
\end{lstlisting}

    \[
        \prod_{i=1}^{n}x_i
    \]

    \begin{center}
        \texttt{\textbackslash sum} $\rightarrow$ soma
        \qquad
        \texttt{\textbackslash prod} $\rightarrow$ produto
    \end{center}

\end{frame}

% =================================================
% NOTAÇÃO AVANÇADA
% =================================================

\section{Derivadas e Integrais}

\begin{frame}[fragile]{Derivadas}

    Notação simples:

\begin{lstlisting}
f'(x)
\end{lstlisting}

    \[
        f'(x)
    \]

    Derivada como fração:

\begin{lstlisting}
\frac{df}{dx}
\end{lstlisting}

    \[
        \frac{df}{dx}
    \]

\end{frame}

\begin{frame}[fragile]{Derivadas de ordem superior}

    Primeira derivada:

\begin{lstlisting}
\frac{df}{dx}
\end{lstlisting}

    \[
        \frac{df}{dx}
    \]

    Segunda derivada:

\begin{lstlisting}
\frac{d^2f}{dx^2}
\end{lstlisting}

    \[
        \frac{d^2f}{dx^2}
    \]

\end{frame}

\begin{frame}[fragile]{Integrais}

    Integral indefinida:

\begin{lstlisting}
\int f(x)\,dx
\end{lstlisting}

    \[
        \int f(x)\,dx
    \]

    Integral definida:

\begin{lstlisting}
\int_a^b f(x)\,dx
\end{lstlisting}

    \[
        \int_a^b f(x)\,dx
    \]

\end{frame}

\begin{frame}[fragile]{Integrais}

    Derivada Parcial:

\begin{lstlisting}
\frac{\parial f}{\partial x}
\end{lstlisting}

    \[
        \frac{\partial f}{\partial x}
    \]

    Gradiente:

\begin{lstlisting}
\nabla f(x,y)
\end{lstlisting}

    \[
    \nabla f(x,y) = \left( \frac{\partial f}{\partial x}, \frac{\partial f}{\partial y} \right)
    \]

\end{frame}

\begin{frame}[fragile]{Integrais}

    
    Integrais Múltiplas:

\begin{lstlisting}
\iint_D f(x,y)
\end{lstlisting}

    \[
    \iint_D f(x,y) \,dx\,dy
    \]

    Caminho/Superfície:

\begin{lstlisting}
\oint_C \mathbf{F} \cdot d\mathbf{r}
\end{lstlisting}

    \[
    \oint_C \mathbf{F} \cdot d\mathbf{r} = \iint_S (\nabla \times \mathbf{F}) \cdot d\mathbf{S}
    \]

\end{frame}

\section{Notação Avançada}

\begin{frame}[fragile]{Conjuntos}

    Alguns símbolos importantes:

\begin{lstlisting}
x \in A
A \subseteq B
A \cup B
A \cap B
\end{lstlisting}

    \[
        x\in A
        \qquad
        A\subseteq B
    \]

    \[
        A\cup B
        \qquad
        A\cap B
    \]

\end{frame}

\begin{frame}[fragile]{Lógica matemática}

    Comandos frequentes:

\begin{lstlisting}
\forall
\exists
\neg
\land
\lor
\Rightarrow
\Leftrightarrow
\end{lstlisting}

    \[
        \forall x,\quad
        P(x)\Rightarrow Q(x)
    \]

    \[
        P(x)\Leftrightarrow Q(x)
    \]

\end{frame}

\begin{frame}[fragile]{Notação de conjuntos}

    Uma expressão mais elaborada:

\begin{lstlisting}
\left\{
x\in\mathbb{R}
\mid x>0
\right\}
\end{lstlisting}

    \[
        \left\{
            x\in\mathbb{R}
            \mid x>0
        \right\}
    \]

    \begin{center}
        Delimitadores + símbolos + condição.
    \end{center}

\end{frame}

\begin{frame}[fragile]{Probabilidade e estatística}

    Algumas notações comuns:

\begin{lstlisting}
P(A\mid B)
E[X]
P(A\cap B)
\end{lstlisting}

    \[
        P(A\mid B)
        \qquad
        E[X]
    \]

    \[
        P(A\cap B)
        =
        P(A\mid B)P(B)
    \]

\end{frame}

\begin{frame}[fragile]{Notação Assintótica}

\begin{lstlisting}
T(n) = \mathcal{O}(n \log n)
\Omega(n) \le T(n) \le \mathcal{O}(n^2)
T(n) = \Theta(n^2)
\end{lstlisting}

    \[
    T(n) = \mathcal{O}(n \log n)
    \]
    \[
    \Omega(n) \le T(n) \le \mathcal{O}(n^2)
    \]
    \[
    T(n) = \Theta(n^2)
    \]

\end{frame}

\begin{frame}[fragile]{Normas e notação vetorial}

    Norma:

\begin{lstlisting}
\|x\|
\end{lstlisting}

    \[
        \|x\|
    \]

    Exemplo:

\begin{lstlisting}
\left\{
x\in\mathbb{R}^n
\mid
\|x\|=1
\right\}
\end{lstlisting}

    \[
        \left\{
            x\in\mathbb{R}^n
            \mid
            \|x\|=1
        \right\}
    \]

\end{frame}

% =================================================
% TEOREMAS E DEMONSTRAÇÕES
% =================================================

\section{Ambientes de Teorema e Demonstração}

\begin{frame}[fragile]{Teoremas e Demonstrações}

Preâmbulo:
\begin{lstlisting}
\usepackage{amsthm}
\newtheorem{theorem}{Teorema}
\newtheorem{lemma}{Lema}
\end{lstlisting}

Corpo:
\begin{lstlisting}
\begin{theorem}[nome do teorema]
    ...
\end{theorem}

\begin{proof}
    ...
\end{proof}
\end{lstlisting}

\end{frame}

\begin{frame}[fragile]{Teoremas e Demonstrações: Ida e Volta}

\begin{lstlisting}
\medskip
\noindent (\textbf{$\Rightarrow$}) 
\medskip
\noindent (\textbf{$\Leftarrow$})
\end{lstlisting}

    \begin{proof}
    Queremos provar que o operador $T$ é contínuo se e somente se é limitado.

    \medskip
    \noindent (\textbf{$\Rightarrow$}) Assume-se que $T$ é contínuo. Então, para todo $\epsilon > 0$, existe $\delta > 0$ tal que...

    \medskip
    \noindent (\textbf{$\Leftarrow$}) Assume-se que $T$ é limitado. Existe uma constante $M > 0$ tal que $\|Tx\| \le M\|x\|$...
    \end{proof}

\end{frame}

\begin{frame}[fragile]{Teoremas e Demonstrações: Prova por Casos}

\begin{lstlisting}
\begin{itemize}
\item[\textbf{Caso x:}] \textit{...}
\end{itemize}
\end{lstlisting}

    \begin{proof}
    Dividimos a análise em dois casos exaustivos:
    \begin{itemize}
        \item[\textbf{Caso 1:}] \textit{$n$ é par ($n = 2k$).} \\
        $n^2 + n = 4k^2 + 2k = 2(2k^2 + k)$, que é par.
        
        \item[\textbf{Caso 2:}] \textit{$n$ é ímpar ($n = 2k + 1$).} \\
        $n^2 + n = (2k+1)^2 + (2k+1) = 4k^2 + 6k + 2 = 2(2k^2 + 3k + 1)$, que é par.
    \end{itemize}
    Logo, $n^2 + n$ é sempre par.
    \end{proof}

\end{frame}

\begin{frame}[fragile]{Teoremas e Demonstrações: Indução Matemática}

\begin{lstlisting}
\begin{description}
\item[Base / Hipotese / Passo :] 
\end{description}
\end{lstlisting}

   \begin{proof}
    Provamos que $\sum_{i=1}^{n} (2i - 1) = n^2$ para todo $n \ge 1$:

    \begin{description}\setlength{\itemsep}{0pt}
        \item[Base:] Para $n=1 \implies 2(1)-1 = 1 = 1^2$. \checkmark
        \item[Hipótese:] Assuma $\sum_{i=1}^{k} (2i - 1) = k^2$ para $k \ge 1$.
        \item[Passo:] Para $n=k+1$, temos:
    \end{description}
    \[
    \sum_{i=1}^{k+1} (2i - 1) = \underbrace{\sum_{i=1}^{k} (2i - 1)}_{k^2} + (2k + 1) = k^2 + 2k + 1 = (k+1)^2 \quad \square
    \]
    \end{proof}

\end{frame}

\begin{frame}[fragile]{Teoremas e Demonstrações: Algébrico com Justificativas}

\begin{lstlisting}
\begin{align*}
\end{align*}
\end{lstlisting}

   \begin{proof}
    Desenvolvendo a expressão, temos:
    \begin{align*}
    (x+1)^3 - (x-1)^3 &= (x^3 + 3x^2 + 3x + 1) - (x^3 - 3x^2 + 3x - 1) && \text{(Expansão dos cubos)} \\
                    &= x^3 - x^3 + 3x^2 + 3x^2 + 3x - 3x + 1 + 1 && \text{(Reagrupamento)} \\
                    &= 6x^2 + 2 && \text{(Simplificação final)}
    \end{align*}
    O que conclui a demonstração.
    \end{proof}

\end{frame}

% =================================================
% ORGANIZAÇÃO
% =================================================

\section{Organização de Equações}

\begin{frame}[fragile]{Expressões em várias linhas}

\begin{lstlisting}
\begin{align}
f(x)
&= x^2 + 2x + 1\\
&= (x+1)^2
\end{align}
\end{lstlisting}

    \[
        \begin{aligned}
            f(x)
            &=x^2+2x+1\\
            &=(x+1)^2
        \end{aligned}
    \]

    \begin{alertblock}{Observe}
        O alinhamento acontece no símbolo marcado por \texttt{\&}.
    \end{alertblock}

\end{frame}

\begin{frame}[fragile]{Com ou sem numeração?}

    \textbf{Com numeração:}

\begin{lstlisting}
\begin{align}
x+y &= 10\\
x-y &= 2
\end{align}
\end{lstlisting}

    \textbf{Sem numeração:}

\begin{lstlisting}
\begin{align*}
x+y &= 10\\
x-y &= 2
\end{align*}
\end{lstlisting}

    \begin{center}
        O \texttt{*} normalmente indica a versão sem numeração.
    \end{center}

\end{frame}

% =================================================
% BOAS PRÁTICAS
% =================================================

\section{Boas Práticas}

\begin{frame}[fragile]{Escreva a estrutura, não a aparência}

    \textbf{Evite:}

\begin{lstlisting}
$$
sin(x)=1
$$
\end{lstlisting}

    \textbf{Prefira:}

\begin{lstlisting}
\[
\sin(x)=1
\]
\end{lstlisting}

    \[
        \sin(x)=1
    \]

\end{frame}

\begin{frame}[fragile]{Agrupamento é fundamental}

    Observe a diferença:

\begin{lstlisting}
x^n+1
\end{lstlisting}

    \[
        x^n+1
    \]

\begin{lstlisting}
x^{n+1}
\end{lstlisting}

    \[
        x^{n+1}
    \]

    \begin{alertblock}{Regra}
        Use \texttt{\{ \}} quando uma estrutura possui mais de um caractere.
    \end{alertblock}

\end{frame}

\begin{frame}{Pensando em estruturas}

    \begin{center}

        \Large
        \textbf{Expressão matemática}

        \vspace{0.5cm}

        $\Downarrow$

        \vspace{0.5cm}

        \textbf{Estruturas menores}

        \vspace{0.5cm}

        $\Downarrow$

        \vspace{0.5cm}

        \textbf{Comandos + argumentos + agrupamento}

    \end{center}

    \vspace{0.5cm}

    \begin{center}
        Uma expressão complexa é construída pela combinação
        de estruturas simples.
    \end{center}

\end{frame}

% =================================================
% DESAFIOS
% =================================================

\section{Exercícios de Transição}

\begin{frame}{Agora é a sua vez}

    \begin{center}

        \Huge
        \textbf{Vamos praticar!}

        \vspace{1cm}

        \Large

        Das estruturas simples

        \[
            x^2,\qquad
            \frac{a}{b},\qquad
            \sqrt{x}
        \]

        às expressões complexas.

    \end{center}

\end{frame}

\begin{frame}{Desafio 1}

    Reproduza em LaTeX:

    \[
        \frac{x^2+3x+2}{x-1}
    \]

    \vspace{0.5cm}

    \begin{center}
        \textbf{Quais estruturas você precisa combinar?}
    \end{center}

\end{frame}

\begin{frame}{Desafio 2}

    Reproduza:

    \[
        \sum_{i=1}^{n}i^2
    \]

    \vspace{0.5cm}

    \begin{center}
        Pense na estrutura antes de pensar no comando.
    \end{center}

\end{frame}

\begin{frame}{Desafio 3}

    Reproduza:

    \[
        A=
        \begin{pmatrix}
            a_{11} & a_{12}\\
            a_{21} & a_{22}
        \end{pmatrix}
    \]

    \vspace{0.4cm}

    \begin{center}
        Matriz + índices + alinhamento.
    \end{center}

\end{frame}

\begin{frame}{Desafio 4}

    Reproduza:

    \[
        \begin{cases}
            x+y=10\\
            2x-y=3
        \end{cases}
    \]

    \vspace{0.4cm}

    \begin{center}
        Qual ambiente representa melhor essa estrutura?
    \end{center}

\end{frame}

\begin{frame}{Desafio 5}

    Reproduza:

    \[
        f(x)=
        \frac{
            \displaystyle\sum_{i=1}^{n}x_i^2
        }{
            \sqrt{x^2+1}
        }
    \]

    \vspace{0.3cm}

    \begin{center}
        \textbf{Não tente decorar.}\\
        Identifique cada estrutura e monte a expressão.
    \end{center}

\end{frame}

\begin{frame}{Desafio 6}

    Reproduza:

    \[
        \nabla f(x,y,z) = \left( \frac{\partial f}{\partial x}, \frac{\partial f}{\partial y}, \frac{\partial f}{\partial z} \right)
    \]

    \vspace{0.3cm}

%    \begin{center}
%        \textbf{Não tente decorar.}\\
%        Identifique cada estrutura e monte a expressão.
%    \end{center}

\end{frame}

\begin{frame}{Desafio 7}

    Reproduza:

    \[
        \nabla f(x) = A^T(Ax - b) =
        \begin{pmatrix}
            0 \\
            0
        \end{pmatrix}
    \]

    \vspace{0.3cm}

%    \begin{center}
%        \textbf{Não tente decorar.}\\
%        Identifique cada estrutura e monte a expressão.
%    \end{center}

\end{frame}

\begin{frame}{Desafio 8}

    Reproduza:

    \[
        H =
        \begin{bmatrix}
            \frac{\partial^2 f}{\partial x^2} & \frac{\partial^2 f}{\partial x \partial y} \\
            \frac{\partial^2 f}{\partial y \partial x} & \frac{\partial^2 f}{\partial y^2}
        \end{bmatrix}
    \]

    \vspace{0.3cm}

    \begin{center}
        \textbf{Não tente decorar.}\\
        Identifique cada estrutura e monte a expressão.
    \end{center}

\end{frame}

\begin{frame}{Desafio 9}

    Reproduza:

A complexidade temporal do algoritmo no pior caso é dada pelo somatório:
\[
T(n) = \sum_{i=1}^{n} \left( \frac{n}{i} \right) = n \cdot H_n
\]

Como a série harmónica satisfaz $H_n = \ln n + \gamma + \mathcal{O}\left(\frac{1}{n}\right)$, a complexidade assintótica é:
\[
T(n) = \mathcal{O}(n \log n)
\]

    \vspace{0.3cm}

%    \begin{center}
%        \textbf{Não tente decorar.}\\
%        Identifique cada estrutura e monte a expressão.
%    \end{center}

\end{frame}

\begin{frame}[fragile]{Último Desafio}

\begin{theorem}[Teorema 1]
Se $f(x) = x^3$, então a área sob a curva $f'(x)$ no intervalo $[0, 2]$ é $8$.
\end{theorem}

\begin{proof}
A derivada de $f(x)$ é $f'(x) = 3x^2$. Calculando a integral definida:
\[
\int_{0}^{2} 3x^2 \, dx = \left[ x^3 \right]_{0}^{2}
\]
Avaliando nos limites de integração pelo Teorema Fundamental do Cálculo:
\[
\left[ x^3 \right]_{0}^{2} = 2^3 - 0^3 = 8
\]
\end{proof}

\end{frame}

% =================================================
% SÍNTESE
% =================================================

\section{Síntese}

\begin{frame}{O que aprendemos?}

    \begin{columns}

        \column{0.48\textwidth}

        \textbf{Sintaxe}

        \[
            x^2
            \qquad
            x_i
            \qquad
            \frac{a}{b}
        \]

        \textbf{Funções}

        \[
            \sin,\quad
            \log,\quad
            \ln
        \]

        \textbf{Notação Assintótica}

        \[
            T(n),
            \quad
            \Omega(n),
            \quad
            \mathcal{O}
        \]

        \column{0.48\textwidth}

        \textbf{Estruturas}

        \[
            \sum,\quad
            \int,\quad
            \lim
        \]

        \textbf{Ambientes}

        \[
            \text{matrix},
            \quad
            \text{align},
            \quad
            \text{cases}
        \]

        \textbf{Demonstrações}

        \[
            \text{theorem},
            \quad
            \text{proof}
        \]

    \end{columns}

\end{frame}

\begin{frame}{A regra mais importante}

    \begin{center}

        \Large
        \textbf{Pense na estrutura, não na aparência.}

        \vspace{1cm}

        \[
            \boxed{
            \text{estrutura}
            +
            \text{estrutura}
            +
            \text{estrutura}
            =
            \text{expressão}
            }
        \]

        \vspace{0.8cm}

        \normalsize

        O LaTeX fica muito mais fácil quando aprendemos
        a reconhecer esses padrões.

    \end{center}

\end{frame}

% =================================================

\end{document}
```
